\documentclass[10pt]{article}
\usepackage[T1]{fontenc}
\usepackage{lmodern,amsmath,amssymb,amsthm,microtype}
\usepackage[margin=0.9in]{geometry}
\usepackage[colorlinks=true,urlcolor=blue,citecolor=blue]{hyperref}
\newtheorem{theorem}{Theorem}
\newtheorem{lemma}{Lemma}
\newtheorem{corollary}{Corollary}
\newcommand{\Pz}{P^{\!*}}
\title{Nontrivial Multiples Preserving the Nonzero Digit Product\\and an Exact Finite Formula for OEIS A087304}
\author{Qichao Wang\\{\small Hebei University of Technology}\\[0.3em]
Daoyu Dong\\{\small University of Electronic Science and Technology of China}}
\date{October 26, 2026}
\hypersetup{pdftitle={Nontrivial Multiples Preserving the Nonzero Digit Product},pdfauthor={Qichao Wang and Daoyu Dong}}
\begin{document}
\maketitle
\begin{abstract}
Every positive integer has infinitely many multiples, other than shifts
of its own digit string by trailing zeros, with the same product of
nonzero digits. A repunit construction proves the nonvanishing conjecture
recorded in OEIS A087304, in every integer base. For the decimal sequence
we give an exact finite digit recurrence for its minimum. Its bound
$n\tau(P)(d+2)$ depends only on the index $n$, its nonzero digit product
$P$, and its digit length $d$.
\end{abstract}

\section{A constructive existence proof}
Let $\Pz_b(n)$ be the product of the nonzero base-$b$ digits of $n$.
The empty product is $1$. OEIS A087304 \cite{oeis} asks for the smallest
multiple $M$ of $n$ satisfying
\[
 \Pz_{10}(M)=\Pz_{10}(n),\qquad
 M\notin\{n10^j:j\ge0\}.
\]
Its entry records the conjecture that this minimum is always positive.
The related A340204 \cite{related}, which includes zeros in its digit
product, already records the usual repunit construction for numbers
coprime to $10$. We use that elementary idea after separating the
base-supported part of the modulus.

\begin{lemma}\label{lem:repunit}
If $\gcd(m,b)=1$, there is a $t\in\{1,\ldots,m\}$ such that
$m\mid R_{b,t}$, where $R_{b,t}=1+b+\cdots+b^{t-1}$.
\end{lemma}
\begin{proof}
Consider the $m$ residues of $R_{b,1},\ldots,R_{b,m}$. If none is zero,
two coincide. Their difference is $b^jR_{b,t}$ for some $1\le t<m$,
and $b$ is invertible modulo $m$. This also includes $m=1$.
\end{proof}

\begin{theorem}\label{thm:exist}
Fix $b\ge2$ and $n\ge1$. Factor $n=qm$ with $\gcd(m,b)=1$ and every
prime divisor of $q$ dividing $b$. Choose $s\ge0$ with $q\mid b^s$,
and choose $t\le m$ as in Lemma~\ref{lem:repunit}. Then, for every $j\ge1$,
\begin{equation}\label{eq:witness}
 M_j=(nb^{jt}+R_{b,jt})b^s
\end{equation}
is divisible by $n$, preserves $\Pz_b(n)$, and is not $nb^h$ for any
$h\ge0$. These witnesses are unbounded. In particular, A087304 has no
zero terms, and a witness has at most $d_b(n)+s+m$ digits.
\end{theorem}
\begin{proof}
The identity
$R_{b,jt}=R_{b,t}(1+b^t+\cdots+b^{(j-1)t})$ shows that
$m\mid R_{b,jt}$. Hence $n\mid b^sR_{b,jt}$ and $n\mid M_j$.
The representation of $M_j$ is precisely the digits of $n$, followed
by $jt$ ones and $s$ zeros, so its nonzero digit product is unchanged.
Finally,
\[
 b^{jt+s}<M_j/n=b^s(b^{jt}+R_{b,jt}/n)<b^{jt+s+1},
\]
using $0<R_{b,jt}<b^{jt}$ and $2\le b$. No integral power of $b$ lies
strictly between these consecutive powers. Taking $j=1$ gives the
digit bound, and letting $j$ increase gives infinitely many witnesses.
\end{proof}

\section{An exact finite formula for the least multiple}
We now work in base ten. Put $P=\Pz_{10}(n)$ and $d=d_{10}(n)$, and write
the digits of $n$ as $c_1\cdots c_d$. Let $\mathcal D(P)$ be the positive
divisors of $P$, and put
\[
 \mathcal S=\{0,\ldots,n-1\}\times\mathcal D(P)\times\{0,\ldots,d+1\},
 \qquad B=|\mathcal S|=n\tau(P)(d+2).
\]
A state $(r,q,h)$ records the current remainder, the unused factor of
the target product, and a state recognizing the forbidden words
$c_1\cdots c_d0^*$. Here $h<d$ means that exactly $h$ initial digits
have matched, $h=d$ means a complete match followed only by zeros,
and $h=d+1$ means a permanent mismatch. Define
\[
 H(h,a)=
 \begin{cases}
 h+1,&h<d\text{ and }a=c_{h+1},\\
 d+1,&h<d\text{ and }a\ne c_{h+1},\\
 d,&h=d\text{ and }a=0,\\
 d+1,&\text{otherwise}.
 \end{cases}
\]
For $a\in\{0,\ldots,9\}$ let $w(a)=\max(1,a)$. A digit is admissible
when $w(a)\mid q$, and the transition is
\begin{equation}\label{eq:transition}
 T_a(r,q,h)=((10r+a)\bmod n,\ q/w(a),\ H(h,a)).
\end{equation}
Introduce a separate initial state $*$, with transitions corresponding
to the admissible digits $1,\ldots,9$ applied to $(0,P,0)$. Thus leading
zeros and the empty word are excluded. The accepting states are
\[
 \mathcal F=\{(0,1,h):h\ne d\}.
\]
Define $D_0(*)=0$, all other $D_0=+\infty$, and for every state $v$ set
\begin{equation}\label{eq:dp}
 D_{\ell+1}(v)=\min_{u\xrightarrow{\ a\ }v}
       \bigl(10D_\ell(u)+a\bigr),\qquad \min\varnothing=+\infty.
\end{equation}

\begin{theorem}[All-index minimum formula]\label{thm:min}
With the explicit states and transitions above,
\begin{equation}\label{eq:minimum}
 \mathrm{A087304}(n)=\min_{1\le\ell\le B}\ \min_{v\in\mathcal F}D_\ell(v).
\end{equation}
In particular its least witness has at most $B$ decimal digits.
\end{theorem}
\begin{proof}
Induction on the word length proves the interpretations of $r,q,h$.
Consequently the accepted words are exactly the eligible multiples of
$n$. Their positivity and divisibility exclude values below $n$, and
the forbidden-word state excludes $n$ and all its decimal shifts.
Theorem~\ref{thm:exist} ensures at least one accepted word exists.

On a shortest accepting path, no state can repeat: removing the segment
between two occurrences leaves an accepting path. If the repetition is
after the initial state, its first nonzero digit is retained; the separate
initial state has no incoming edges. There are $B+1$ states including
$*$, so such a path has at most $B$ edges. A positive decimal integer
with fewer digits is smaller than one with more digits. For each fixed
length, \eqref{eq:dp} gives the smallest label reaching each state:
appending the same digit preserves the ordering of prefix values.
Thus \eqref{eq:minimum} gives the global least eligible integer.
\end{proof}

Formula~\eqref{eq:minimum} is a finite minimum with a fixed number of digit
recurrence steps; it does not require scanning successive multiples.
The accompanying program evaluates it by breadth-first search with
digits ordered increasingly. State merging preserves the shortest,
lexicographically first prefix. Independent direct enumeration verifies
the initial published minima. The already recorded identity
$a(10n)=10a(n)$, credited to Chai Wah Wu (2017) in \cite{oeis}, also follows
by deleting or appending one trailing zero in an eligible multiple.

\begin{thebibliography}{9}
\bibitem{oeis} A. Murthy, OEIS A087304,
\url{https://oeis.org/A087304}, accessed October 26, 2026.
\bibitem{related} B. Schott, OEIS A340204,
\url{https://oeis.org/A340204}, accessed October 26, 2026.
\end{thebibliography}
\end{document}
